% TOP_PROBLEM: {"id": 30002744, "title": "Reflection-Positivity Research Questions", "category_id": 16, "category": {"id": 16, "name": "physics", "display_name": "Mathematical Physics", "description": "Problems at the intersection of mathematics and physics.", "slug": "physics", "order_index": 16, "created_at": "2026-07-31T15:26:25.671Z"}, "status": "open", "statusQualification": "This is the report\\textquotesingle s bundled problem list. All requested connections, realizations and removal of growth assumptions are preserved below; the common-core question is programmatic and no unlisted set of axioms is supplied as an alleged original theorem. The full primary list also asks for a Lax--Phillips-type normal form for one-parameter groups whose positive part preserves a standard real subspace, reflection positivity in situations without a Euclidean Hilbert space, and duality for multiple reflection-positive involutions. It asks about multiparticle cluster-expansion spectra beyond a mass gap, and OS--Wightman equivalence without growth estimates or regularity assumptions. These specific directions are retained as additional source context to the condensed catalogue wording. The report explicitly distinguishes deep open problems from background questions in its list."}
% FIRST_POSTED: 2026-10-11
% ENTRY_KIND: statement_only
% SUBMISSION: {"kind": "statement_only", "category": "top_problems", "problem_id": "30002744", "proposal_type": "established", "source": "Mini-Workshop: Reflection Positivity in Representation Theory, Stochastics and Physics (2014), Oberwolfach Reports, DOI 10.4171/owr/2014/55, pp. 3059--3102. Primary report accessed 2026-10-11; relevant PDF pages 39, 40. https://doi.org/10.4171/owr/2014/55", "status_evidence": "Canonical source status open retained. Supplied literature triage dated 2026-08-21: The record is a list of twenty heterogeneous reflection-positivity questions, some conceptual and some deep; no single later result resolves the list as a whole. Primary-report access checked 2026-10-11; this historical report and inherited triage do not certify current resolution of every component.", "status_checked_on": "2026-10-11", "duplicate_check": "Compared canonical ID, title and normalized target with the complete local UnsolvedMath registry; normalized exact-text matches: none; title matches: none. Related or overlapping canonical records: none identified in this batch. Primary source context and the existing OWR statement audit used for scope and duplicate review; an exhaustive mathematical-equivalence/literature search was not performed.", "attribution": "Original proposer not named in the supplied canonical record; attribution remains to the identified Oberwolfach contribution/report, not its editors as invented proposers. Explanatory statement preparation and source/semantic audit assisted by OpenAI GPT-6 Codex; no human mathematical authorship or independent proof verification is claimed."}
% SOURCE_STATEMENT: "Resolve the following questions about reflection positivity: describe super versions of massive free fields and reflection-positive Hilbert spaces; determine the role of antiunitary involutions; construct Euclidean realizations of Borchers triples and algebraic quantum field theories; find diffeomorphism-invariant examples and a noncommutative-geometric theory; clarify links with complex measures, stochastic quantization, number theory, string theory, one-dimensional conformal field theory, cluster expansions, modular theory, and the Bessis–Moussa–Villani theorem; remove growth assumptions from the equivalence of the Osterwalder–Schrader and Wightman axioms; and characterize the common structural core of these phenomena."
% Preparation and final self-audit: CONTRIBUTING.md prompts C and D; source/audit artifacts retained in the local physics_statements_2026_10_11 batch.
% Imported from: physics_statements.tex; UnsolvedMath ID 30002744
% !TeX program = lualatex
\documentclass[11pt]{article}
\usepackage[margin=25mm]{geometry}
\usepackage{fontspec}
\setmainfont{Latin Modern Roman}
\usepackage{amsmath,amssymb,mathtools}
\usepackage{unicode-math}
\setmathfont{Latin Modern Math}
\usepackage{xurl,hyperref}
\hypersetup{colorlinks=true,urlcolor=blue}
\setcounter{secnumdepth}{0}
\setlength{\parindent}{0pt}
\setlength{\parskip}{5pt}
\setlength{\emergencystretch}{3em}
\newcommand{\sref}[2]{\hyperlink{source:#1:#2}{[S#2]}}
\newcommand{\cataloguescope}[1]{\paragraph{Recorded target.}#1}
\newcommand{\R}{\mathbb R}
\newcommand{\C}{\mathbb C}
\newcommand{\N}{\mathbb N}
\newcommand{\Z}{\mathbb Z}
\newcommand{\Q}{\mathbb Q}
\newcommand{\F}{\mathbb F}
\newcommand{\eps}{\varepsilon}
\begin{document}
% UnsolvedMath ID 30002744; source catalogue code OWR-13357-003
\setcounter{section}{30002743}
\section{Reflection-Positivity Research Questions}\label{30002744}
{\small UnsolvedMath ID 30002744 · Mathematical Physics}
\subsection{Definitions and mathematical statement}

\paragraph{Shared notation.}$\mathbb N=\{1,2,\ldots\}$; the standard number fields have their usual meanings. All additional conventions are stated locally.


\paragraph{Definitions and conventions.}
In its basic form, reflection positivity concerns a Hilbert space $\mathcal E$, a unitary involution $\theta$ and a subspace $\mathcal E_+$ such that $\langle\theta v,v\rangle\geq0$ for $v\in\mathcal E_+$. Quotienting by the null space produces a physical Hilbert space. Osterwalder--Schrader axioms describe Euclidean correlation functions; Wightman axioms describe relativistic fields. A Borchers triple consists of an operator algebra, translation representation and vacuum with compatible wedge-locality properties. Super versions incorporate a $\mathbb Z_2$ grading.

\paragraph{Target retained from the canonical source catalogue.}
\paragraph{Catalogue wording retained for traceability.}
Resolve the following questions about reflection positivity: describe super versions of massive free fields and reflection-positive Hilbert spaces; determine the role of antiunitary involutions; construct Euclidean realizations of Borchers triples and algebraic quantum field theories; find diffeomorphism-invariant examples and a noncommutative-geometric theory; clarify links with complex measures, stochastic quantization, number theory, string theory, one-dimensional conformal field theory, cluster expansions, modular theory, and the Bessis–Moussa–Villani theorem; remove growth assumptions from the equivalence of the Osterwalder–Schrader and Wightman axioms; and characterize the common structural core of these phenomena.

\paragraph{Source scope and qualification.}
This is the report\textquotesingle s bundled problem list. All requested connections, realizations and removal of growth assumptions are preserved below; the common-core question is programmatic and no unlisted set of axioms is supplied as an alleged original theorem. The full primary list also asks for a Lax--Phillips-type normal form for one-parameter groups whose positive part preserves a standard real subspace, reflection positivity in situations without a Euclidean Hilbert space, and duality for multiple reflection-positive involutions. It asks about multiparticle cluster-expansion spectra beyond a mass gap, and OS--Wightman equivalence without growth estimates or regularity assumptions. These specific directions are retained as additional source context to the condensed catalogue wording. The report explicitly distinguishes deep open problems from background questions in its list.
\subsection{Short English statement}
Explain and extend reflection positivity across its field-theoretic, algebraic, geometric and probabilistic appearances.
\subsection{Sources}
\begin{itemize}
\item[\hypertarget{source:30002744:1}{S1}] ULAM.AI and the UnsolvedMath Contributors, UnsolvedMath, record 30002744 (OWR-13357-003). Catalogue curation under CC BY 4.0; definitions and English exposition added locally.
\url{https://huggingface.co/datasets/ulamai/UnsolvedMath}.
\item[\hypertarget{source:30002744:2}{S2}] Mini-Workshop: Reflection Positivity in Representation Theory, Stochastics and Physics (2014), Oberwolfach Reports, DOI 10.4171/owr/2014/55, pp. 3059--3102. Primary report accessed 2026-10-11; relevant PDF pages 39, 40.
\url{https://doi.org/10.4171/owr/2014/55}.
\end{itemize}
\label{end:30002744}
\end{document}
